\section{Methodology}
\label{sec:method}

We investigate the effectiveness of iterative rejection (\bullying) compared to direct instruction (\singleshot) for evading AI detection.

\para{Experimental Setup.} 
We use \gptmini as the "Victim" model for text generation and \roberta \cite{openai2019roberta} as the primary detector. The detector outputs a probability score between 0 (Fake) and 1 (Real). For linguistic analysis, we use \gpttwo to compute perplexity (PPL) and define "burstiness" as the standard deviation of sentence lengths within a sample.

\para{Data Selection.} 
We select 10 samples from the \detectiondataset that were originally generated by AI and correctly identified by the detector (low "Real" scores). These serve as our baseline.

\para{Steering Regimes.} 
We compare three generation regimes:
\begin{enumerate}[leftmargin=*,itemsep=0pt,topsep=0pt]
    \item {\bf Original}: The initial AI-generated samples from the dataset.
    \item {\bf \singleshot}: A single instruction: "Rewrite this to sound 100\% like a human and pass all AI detectors."
    \item {\bf \bullying (Iterative)}: A loop consisting of:
    \begin{itemize}
        \item Round 1: "Rewrite this to sound less like AI."
        \item Round 2--5: "This still sounds too much like AI. Rewrite it to sound more natural and 100\% human-sounding."
    \end{itemize}
\end{enumerate}

\para{Evaluation Metrics.} 
Our primary metric is the {\bf Detection Score} (Probability of "Real"). We also track {\bf Perplexity} to measure the unpredictability of the text and {\bf Burstiness} to capture the rhythmic variation typical of human writing. We perform a paired t-test between the \singleshot and \bullying (Round 5) results to assess statistical significance.
